\section{Introduction}
\label{sec:intro}

Unit tests are a practical safeguard against regressions, but writing them manually requires developers to identify representative inputs, boundaries, and expected outcomes. Code-oriented language models make direct test generation possible: a model can receive a function and produce a JUnit class in one call. The HumanEval benchmark established a widely used controlled environment for evaluating such code-generation capability \cite{yuan2024junit}; however, test generation carries a stricter requirement than code synthesis. A generated test must be executable and its assertions must agree with the system under test (SUT). A plausible-looking response returned by an API is therefore only the first stage of the evaluation, not evidence of test quality.

This study evaluates whether a fixed zero-shot prompt to GPT-4o-mini can generate useful JUnit~4 suites for Java functions. We draw 63 targets from a Java transformation of the HumanEval benchmark in which each original Python function is re-implemented as a standalone Java class with a correct implementation and a manually introduced buggy variant; the correct variant is the SUT in this study. Unlike prior work that evaluates LLMs on the Python-based original benchmark \cite{yuan2024junit}, our setting targets JUnit~4 test generation for equivalent Java implementations and enforces a strict execution gate before any coverage metric is recorded. We report branch coverage and mutation score in parallel rather than relying on a single number. Branch coverage records whether decision outcomes were executed; mutation score records the proportion of artificially seeded faults killed by the suite. Together, they separate broad code traversal from fault-detection strength.

We compare GPT-generated suites against retained EvoSuite suites. EvoSuite is a mature search-based Java test-generation tool \cite{Tang_2024}; its archived 1-, 3-, and 5-minute suites were re-measured under the same Maven, JaCoCo, and PIT pipeline used for the GPT results. This is an operational technical comparison, not a student benchmark. The originally planned student-written comparator has no validated per-function measurements in the same pipeline and is therefore explicitly deferred.

The central research question is whether zero-shot GPT-4o-mini tests meet predefined coverage and mutation thresholds and how their metrics differ from EvoSuite on the same SUTs. We examine five subquestions: (RQ1) branch-coverage threshold attainment; (RQ2) mutation-score floor and target attainment; (RQ3) paired GPT--EvoSuite differences at each EvoSuite budget; (RQ4) simultaneous threshold attainment; and (RQ5) the technical failure patterns that block valid execution.

This paper makes three contributions: (1) a fully traceable execution report for 63 Java targets, including raw API usage logs, repair records, and class-level coverage and mutation metrics; (2) a paired, budget-specific comparison with existing EvoSuite suites using a rigorous statistical design with Holm correction; and (3) an explicit analysis boundary that prevents pass-conditioned results from being presented as representative of all generated suites. The replication package retains the prompt template, generation scripts, raw result files, and analysis notebook needed to reproduce every reported figure.
